\chapter{Trading on-chain}\label{ch:hf:on-chain-trading}

Une grosse opération de swap sur une plateforme décentralisée écarte le prix de son pool de celui des plateformes centralisées. Dans le même bloc, les searchers rivalisent
pour l'arbitrer en sens inverse, et la plus grande part de leurs gains est cédée au constructeur qui inclut leur bundle. Dans l'enchère au premier prix de ce chapitre, deux
searchers concurrents laissent au constructeur 93~\% de la valeur d'une opportunité, dix en laissent 98~\%~; l'arbitrage lui-même, sur un pool de 20~millions de dollars pendant une
semaine simulée, rapporte à ses searchers 4\,300~\$ après le gas. Le fournisseur de liquidité de l'autre côté de ces transactions, couvert sur une plateforme centralisée,
perd 26\,000~\$ au profit de ces searchers, soit la perte par rapport au rééquilibrage (loss-versus-rebalancing) qu'une formule prédit à 10~\% près, et conserve 313\,000~\$ de frais nets de cette perte.

\section{Arbitrage entre plateformes centralisées et chaîne}

Un pool à produit constant (livre~3, chapitre~20) affiche un prix qui ne bouge que lorsque quelqu'un y négocie. Lorsque les plateformes centralisées
bougent, le pool reste périmé jusqu'à ce qu'un arbitragiste le ramène en sens inverse par une transaction~: l'arbitrage CEX-DEX du livre~3, chapitre~22. La transaction optimale
admet une forme fermée. Pour un pool de $x$ tokens et $y$ dollars, avec des frais $1-g$, et un prix centralisé $p$ inférieur au prix du pool de plus que les frais,
l'arbitragiste vend des tokens dans le pool jusqu'à ce que la sortie marginale du dernier, $gxy/(x+g\,\Delta x)^2$, soit égale à $p$~:
\[
\Delta x^\ast=\frac{\sqrt{gxy/p}-x}{g},\qquad\text{gain}=\frac{g\,\Delta x^\ast y}{x+g\,\Delta x^\ast}-p\,\Delta x^\ast,
\]
et de façon symétrique avec des dollars lorsque le pool est bon marché (\cref{lst:hf:on-chain-trading:arb}). La transaction ramène le pool à moins des frais
près du prix centralisé. La couverture de l'arbitragiste est une transaction sur la plateforme centralisée, qui ne se trouve pas dans le même bloc et n'est pas atomique~: elle porte un risque de prix
entre les deux jambes, ainsi que le risque que le bloc n'inclue pas la jambe on-chain.

\section{Arbitrage atomique et recherche}

\begin{definition}[Arbitrage atomique]\label{def:hf:on-chain-trading:atomic}
Un \emph{arbitrage atomique}\index{arbitrage atomique} est un arbitrage dont toutes les jambes s'exécutent dans une seule transaction on-chain, typiquement sur plusieurs plateformes décentralisées
et souvent avec un flash loan, de sorte que soit chaque jambe réussit, soit la transaction entière est annulée et seuls les frais de la
tentative échouée sont perdus.
\end{definition}

Daian et ses coauteurs ont documenté l'essor des bots d'arbitrage sur les plateformes décentralisées qui, à l'image des traders à haute fréquence, payaient des frais de transaction élevés
et optimisaient leur latence pour passer devant les transactions des utilisateurs ordinaires~; ils ont décrit les enchères de gas prioritaire, dans lesquelles les bots surenchérissent sur les frais pour obtenir des
positions précoces dans le bloc, et ont montré que les frais payés pour l'ordonnancement faisaient peser un risque sur la couche de consensus. Ce risque est ce que la séparation proposeur-constructeur
(livre~3, chapitre~22) a réorganisé~: les searchers envoient désormais leurs bundles aux constructeurs, qui assemblent les blocs, et les proposeurs choisissent le bloc le plus
précieux.

\begin{omfigure}
\begin{tikzpicture}[font=\footnotesize,>=stealth,box/.style={draw,rounded corners,fill=omThm!8,minimum width=2.3cm,minimum height=0.75cm,align=center}]
\node[box] (u) at (0,1.2) {swap d'utilisateur\\(mempool)};
\node[box] (c) at (0,-1.2) {prix des\\plateformes centralisées};
\node[box,fill=omDef!10] (s) at (3.6,0) {les searchers\\($n$ bundles)};
\node[box] (b) at (8.4,0) {constructeur};
\node[box] (p) at (12.9,0) {proposeur\\(bloc)};
\draw[->] (u) -- (s);
\draw[->] (c) -- (s);
\draw[->] (s) -- node[above=3pt,font=\scriptsize] {bundle, enchère} (b);
\draw[->] (b) -- node[below=3pt,font=\scriptsize] {bloc, paiement} (p);
\node[font=\scriptsize,align=center] at (6.0,-1.1) {enchère au premier prix~:\\la meilleure enchère est incluse};
\end{tikzpicture}

\omcaption{Le parcours d'un arbitrage on-chain~: les searchers voient une transaction d'un utilisateur ou un mouvement du prix centralisé, construisent un bundle et enchérissent pour son inclusion dans
une enchère au premier prix menée par un constructeur, qui assemble un bloc et verse au proposeur le montant de son inclusion. Source~: livre~3, chapitre~22~; le présent
chapitre.}\label{fig:hf:on-chain-trading:pipeline}
\end{omfigure}

\section{Bundles, constructeurs et enchères}

\begin{definition}[Enchère de bundle]\label{def:hf:on-chain-trading:bid}
Une \emph{enchère de bundle}\index{enchère de bundle} (bundle bid) est le paiement qu'un searcher attache à un bundle de transactions pour son inclusion dans un bloc, versé au constructeur
(directement ou par les frais de priorité de la transaction) seulement si le bundle est inclus.
\end{definition}

Lorsque de nombreux searchers voient la même opportunité, l'enchère sur les bundles décide qui l'obtient et à quel prix
(\cref{fig:hf:on-chain-trading:pipeline}). Le modèle du chapitre (\cref{lst:hf:on-chain-trading:auction})~: une opportunité vaut 100~\$ sur la base d'une information
publique~; la valeur de chaque searcher est inférieure de son propre coût (gas, inventaire, risque sur la jambe centralisée), uniformément entre 0~\% et 20~\%~; les enchères
sont scellées et la meilleure est payée. Avec des valeurs privées uniformes, l'enchère symétrique d'équilibre minore la valeur vers la plus basse de $1/n$. Un
searcher seul ne paie rien~: c'est un monopoleur. Avec deux, le constructeur reçoit 93~\% de la valeur du gagnant~; avec cinq, 97~\%~; avec vingt, 99~\%
(\cref{fig:hf:on-chain-trading:auction}). Ce que garde le searcher est son avantage de coût sur le deuxième meilleur, et non l'opportunité.

\begin{omfigure}
\begin{tikzpicture}
\begin{axis}[width=11cm,height=5.2cm,axis y line*=right,axis x line=none,xmode=log,log basis x=10,xmin=0.9,xmax=22,ymin=0,ymax=100,
  ylabel={\small le gagnant garde (\$ sur 100~\$)},tick label style={font=\footnotesize},table/col sep=comma]
\addplot[omDef,thick,dashed,mark=square*,mark size=1.3pt] table[x=n,y=searcher]{figdata/hft/25-on-chain-trading/auction.csv};
\label{plot:hf:auckeep}
\end{axis}
\begin{axis}[width=11cm,height=5.2cm,axis y line*=left,xmode=log,log basis x=10,xmin=0.9,xmax=22,xtick={1,2,3,5,10,20},
  xticklabels={1,2,3,5,10,20},xlabel={\small searchers concurrents},ylabel={\small part du constructeur dans la valeur du gagnant (\%)},ymin=0,
  ymax=105,tick label style={font=\footnotesize},grid=major,grid style={gray!25},table/col sep=comma,
  legend style={font=\scriptsize,at={(0.97,0.4)},anchor=east}]
\addplot[omThm,thick,mark=*,mark size=1.4pt] table[x=n,y=share]{figdata/hft/25-on-chain-trading/auction.csv};
\addlegendentry{part du constructeur (gauche)}
\addlegendimage{/pgfplots/refstyle=plot:hf:auckeep}\addlegendentry{ce que garde le gagnant (droite)}
\end{axis}
\end{tikzpicture}

\omcaption{Une enchère au premier prix sur des bundles pour une opportunité de 100~\$ entre searchers dont les valeurs sont inférieures de leurs propres coûts, uniformes entre 0 et
20~\%~: la part du constructeur dans la valeur du searcher gagnant, et ce que garde le gagnant, en fonction du nombre de searchers concurrents (échelle logarithmique)~;
20\,000 enchères par point. Données~: \texttt{hf\_onchain.auctions}.}\label{fig:hf:on-chain-trading:auction}
\end{omfigure}

\begin{dated}{2026-09}{La loi et l'enchère de bloc}\label{dat:hf:on-chain-trading:law}
En novembre 2025, un juge fédéral de Manhattan a prononcé un procès nul dans l'affaire d'Anton et James Peraire-Bueno, poursuivis pour fraude électronique,
complot en vue de commettre une fraude électronique et complot en vue de blanchir de l'argent, à propos d'une exploitation présumée du logiciel MEV-Boost qui aurait rapporté 25~millions de dollars
en douze secondes~; selon les comptes rendus, le jury n'a pas pu s'accorder sur la manière d'appliquer la loi. Le parquet pourrait rejuger l'affaire.
\end{dated}

\section{Liquidations sur les protocoles de prêt}

Les protocoles de prêt (livre~3, chapitre~23) permettent à quiconque de rembourser la dette d'un emprunteur sous-collatéralisé et de prendre son collatéral avec un bonus.
Liquider 1~million de dollars de dette avec un bonus de 5~\% rapporte 50\,000~\$ moins le gas~; de nombreux searchers surveillent le facteur de santé de chaque position, et
l'opportunité est mise aux enchères exactement comme un arbitrage. Les flash loans permettent à un searcher de liquider sans capital, en empruntant, en remboursant la dette, en vendant
le collatéral et en remboursant le prêt dans une seule transaction atomique.

\section{Liquidité just-in-time et apport de liquidité couvert}

Un fournisseur de liquidité dans le pool prend l'autre côté de chaque arbitrage. Milionis, Moallemi, Roughgarden et Zhang ont décomposé son rendement en une
exposition de marché et une partie microstructurelle~: les frais perçus moins les pertes subies face aux arbitragistes. Pour un pool à produit constant de valeur $V$, avec la
volatilité $\sigma$ du token, la perte par rapport au rééquilibrage s'élève à $\sigma^2V/8$ par unité de temps~; un fournisseur qui couvre l'exposition du pool au token sur une
plateforme centralisée conserve les frais diminués de ce taux.

\begin{definition}[Apport de liquidité couvert]\label{def:hf:on-chain-trading:hedged}
L'\emph{apport de liquidité couvert}\index{apport de liquidité couvert} (hedged liquidity provision) consiste à fournir de la liquidité à un teneur de marché automatisé (AMM) tout en couvrant
l'exposition de la position au cours du token sur une autre plateforme, de sorte que le résultat du fournisseur soit égal aux frais du pool diminués de sa perte face aux arbitragistes, avec
le risque lié au cours du token supprimé.
\end{definition}

Le pool du chapitre (\cref{lst:hf:on-chain-trading:pool})~: 20~millions de dollars avec des frais de 0.3~\%, un token valant 3\,000~\$ avec une volatilité quotidienne de 4~\%, des blocs de 12 secondes
pendant une semaine, 2\,000~\$ de swaps de bruit par bloc dans des directions aléatoires, avec un arbitrage après chacun dès lors qu'il rapporte plus de 5~\$ de gas. Sur trois semaines simulées, la
loss-versus-rebalancing varie de 25\,800~\$ à 29\,500~\$, contre 28\,000~\$ selon la formule~; les arbitragistes reprennent aussi l'impact sur le prix que les noise traders
ont payé, soit 20\,000~\$ à 22\,500~\$ par semaine, qui leur reviennent par construction. Le fournisseur couvert encaisse 339\,000~\$ de frais et son résultat net s'élève à 313\,000~\$.

\begin{omfigure}
\begin{tikzpicture}
\begin{axis}[width=11cm,height=5.2cm,xlabel={\small frais du pool (points de base)},ylabel={\small milliers de \$ sur la semaine},xmode=log,log basis x=10,
  xtick={5,10,30,100},xticklabels={5,10,30,100},xmin=4,xmax=120,ymode=log,log basis y=10,ymin=10,ymax=1500,
  tick label style={font=\footnotesize},grid=major,grid style={gray!25},table/col sep=comma,
  legend style={font=\scriptsize,at={(0.03,0.97)},anchor=north west}]
\addplot[omThm,thick,mark=*,mark size=1.4pt] table[x=fee,y=fees]{figdata/hft/25-on-chain-trading/lp.csv};
\addplot[omDef,thick,dashed,mark=square*,mark size=1.3pt] table[x=fee,y=lvr]{figdata/hft/25-on-chain-trading/lp.csv};
\addplot[gray,thick,dotted,mark=triangle*,mark size=1.3pt] table[x=fee,y=net]{figdata/hft/25-on-chain-trading/lp.csv};
\legend{frais perçus,loss-versus-rebalancing,net couvert}
\end{axis}
\end{tikzpicture}

\omcaption{La semaine du fournisseur de liquidité couvert sur un pool de 20~millions de dollars, en fonction de son palier de frais (échelles logarithmiques)~: frais perçus, loss-versus-rebalancing et
le net~; le volume des noise traders est maintenu fixe, ce qui avantage les frais élevés (le volume réel des noise traders baisse avec les frais). Données~:
\texttt{hf\_onchain.by\_fee}.}\label{fig:hf:on-chain-trading:lp}
\end{omfigure}

La loss-versus-rebalancing ne dépend guère des frais (\cref{fig:hf:on-chain-trading:lp})~: elle est fixée par la volatilité et la taille du pool. Les frais déterminent
le revenu, par l'intermédiaire du volume de bruit que le pool attire~; le modèle maintient ce volume fixe, ce qui est justement l'hypothèse qu'un fournisseur ne doit pas
faire lorsqu'il choisit un palier de frais. La liquidité just-in-time (livre~3, chapitre~20) est le cas extrême de l'apport couvert~: un searcher ajoute
de la liquidité concentrée autour d'un swap de grande taille qu'il a repéré dans le mempool, encaisse les frais de ce swap et retire la liquidité dans le même bloc.

\section{Attaques sandwich et droit}

Une attaque sandwich (livre~3, chapitre~22) achète avant le swap d'une victime et vend après lui, en prenant l'impact sur le prix que la victime était prête à
tolérer. Sur le pool du chapitre, un achat de 100\,000~\$ assorti d'une tolérance de glissement de 1~\% peut être précédé d'un front-run de 50\,900~\$ qui laisse la victime
0.33 token, soit environ 990~\$ de perte, exactement sa tolérance~; le gain brut de l'attaquant est de 700~\$, dont, dans une enchère de bloc concurrentielle, 90~\% reviennent
au constructeur et 10~\$ de gas, il reste 60~\$ (\texttt{firm.sandwich}, livre~3). La perte de la victime est le revenu de l'attaque~; la part revenant au constructeur est
prélevée sur celui-ci. Savoir s'il s'agit de fraude, de manipulation de marché ou d'aucune des deux n'a pas été tranché en justice~: la seule poursuite allée jusqu'à un jury, au sujet
d'une exploitation visant les bots sandwich eux-mêmes, s'est soldée par un mistrial (encadré daté). Ce chapitre ne décrit les attaques sandwich que pour chiffrer ce que
perdent ses victimes~; ce n'est pas une fiche de stratégie.

\section{Fiches de stratégie}

\begin{strategyfile}{arbitrage CEX-DEX}\label{strat:hf:on-chain-trading:cexdex}
\sfield{Qui paie, et pourquoi} Les fournisseurs de liquidité des pools aux prix obsolètes, par la loss-versus-rebalancing.
\sfield{Instruments et plateformes} Pools à produit constant et pools à liquidité concentrée~; les carnets d'ordres des plateformes centralisées.
\sfield{Signal} Le cours du pool face au cours centralisé, au-delà des frais et du gas.
\sfield{Dimensionnement et exécution} La taille en forme fermée~; la jambe on-chain dans un bundle, la jambe centralisée en guise de couverture.
\sfield{Coûts} Le gas, l'\omterm{def:hf:on-chain-trading:bid}{enchère de bundle} (93~\% à 99~\% de la valeur face aux concurrents), le risque et les frais de la jambe centralisée.
\sfield{Comment elle s'éteint} La concurrence dans l'\omterm{def:hf:on-chain-trading:bid}{enchère de bundle}~; dans le modèle, l'arbitrage de la semaine rapporte à ses searchers 4\,300~\$ après déduction du gas, avant les enchères.
\sfield{Horizon, capacité, infrastructure} Un bloc~; un inventaire des deux côtés~; les connexions aux constructeurs.
\sfield{Backtest honnête} États du pool bloc par bloc, le cours centralisé à l'heure de chaque bloc, enchères qui auraient été remportées.
\sfield{Sources} Milionis, Moallemi, Roughgarden et Zhang~; le livre~3, chapitre~22.
\end{strategyfile}

\begin{strategyfile}{Arbitrage atomique DEX-DEX}\label{strat:hf:on-chain-trading:atomic}
\sfield{Qui paie, et pourquoi} Les traders dont les swaps laissent deux pools à des prix différents.
\sfield{Instruments et plateformes} Pools d'une même paire ou cycles de paires, sur les plateformes décentralisées (DEX).
\sfield{Signal} Un cycle dont la sortie dépasse l'entrée après frais.
\sfield{Dimensionnement et exécution} Une transaction atomique, souvent financée par un flash loan~; elle est annulée si le cours a bougé.
\sfield{Coûts} Le gas, y compris pour les tentatives annulées~; l'\omterm{def:hf:on-chain-trading:bid}{enchère de bundle}.
\sfield{Comment elle s'éteint} Les enchères de gas prioritaires documentées par Daian et ses coauteurs, et leur successeur, l'enchère des constructeurs.
\sfield{Horizon, capacité, infrastructure} Un bloc~; mempool et simulation de l'état de la chaîne.
\sfield{Backtest honnête} Simuler par rapport à l'état de la chaîne au bloc, y compris les bundles concurrents.
\sfield{Sources} Daian et al.\ (2019, Flash Boys 2.0).
\end{strategyfile}

\begin{strategyfile}{Recherche de liquidations sur protocoles de prêt}\label{strat:hf:on-chain-trading:liquidation}
\sfield{Qui paie, et pourquoi} Les emprunteurs dont le collatéral baisse, par le biais du bonus de liquidation du protocole.
\sfield{Instruments et plateformes} Protocoles de prêt~; les marchés du collatéral où le vendre.
\sfield{Signal} Facteurs de santé inférieurs à un, à partir des positions et des cours on-chain.
\sfield{Dimensionnement et exécution} Rembourser la dette, prendre le collatéral avec le bonus, le vendre~; des flash loans si nécessaire~; enchérir pour être inclus.
\sfield{Coûts} Le gas, l'enchère, la vente du collatéral sur un marché en baisse.
\sfield{Comment elle s'éteint} Enchères encombrées lors des krachs, quand le gas et le glissement du collatéral culminent ensemble.
\sfield{Horizon, capacité, infrastructure} Un bloc~; oracle et suivi des positions.
\sfield{Backtest honnête} Mises à jour de l'oracle telles qu'elles étaient, liquidateurs concurrents, prix de vente lors du krach.
\sfield{Sources} Le livre~3, chapitre~23.
\end{strategyfile}

\begin{strategyfile}{Fourniture de liquidité just-in-time}\label{strat:hf:on-chain-trading:jit}
\sfield{Qui paie, et pourquoi} Les gros swappers, par les frais de leur swap.
\sfield{Instruments et plateformes} Pools à liquidité concentrée.
\sfield{Signal} Un gros swap en attente dans le mempool ou une enchère de flux d'ordres.
\sfield{Dimensionnement et exécution} Ajouter de la liquidité autour du cours pour le swap, la retirer ensuite, en un seul bundle~; couvrir l'inventaire que laisse le swap.
\sfield{Coûts} Le gas, l'\omterm{def:hf:on-chain-trading:bid}{enchère de bundle}, la couverture.
\sfield{Comment elle s'éteint} Le flux d'ordres privé qui masque les swaps jusqu'à leur inclusion.
\sfield{Horizon, capacité, infrastructure} Un bloc.
\sfield{Backtest honnête} Seuls les swaps visibles avant leur inclusion.
\sfield{Sources} Le livre~3, chapitre~20.
\end{strategyfile}

\begin{strategyfile}{Enchérir sur les bundles dans les enchères de constructeurs}\label{strat:hf:on-chain-trading:bidding}
\sfield{Qui paie, et pourquoi} Personne ne paie l'enchérisseur~: l'enchère détermine la part d'une opportunité que le searcher garde.
\sfield{Instruments et plateformes} Enchères de constructeurs pour les bundles.
\sfield{Signal} La valeur de l'opportunité, le coût propre du searcher, le nombre de concurrents.
\sfield{Dimensionnement et exécution} Minorer l'enchère selon l'équilibre~: la rapprocher de la valeur la plus basse de $1/n$, avec des coûts uniformes.
\sfield{Coûts} Enchères perdues~; surenchère.
\sfield{Comment elle s'éteint} Davantage de concurrents~: le gagnant garde 6.7~\$ sur 100~\$ face à un rival et 1.0~\$ face à dix-neuf~; et les exploitations du logiciel
de l'enchère, objet de la poursuite Peraire-Bueno (encadré daté).
\sfield{Horizon, capacité, infrastructure} Un bloc.
\sfield{Backtest honnête} Les enchères des concurrents issues de la chaîne, et non celles du searcher.
\sfield{Sources} Ce chapitre~; l'encadré daté.
\end{strategyfile}

\begin{strategyfile}{Market making couvert à liquidité concentrée}\label{strat:hf:on-chain-trading:hedgedlp}
\sfield{Qui paie, et pourquoi} Les swappers, par les frais, déduction faite de ce que prennent les arbitragistes.
\sfield{Instruments et plateformes} Pools à liquidité concentrée~; une plateforme centralisée pour la couverture.
\sfield{Signal} Revenu des frais comparé au taux de loss-versus-rebalancing $\sigma^2V/8$.
\sfield{Dimensionnement et exécution} Fournir de la liquidité là où les frais dépassent le taux~; couvrir le delta du pool à chaque bloc ou au-delà d'une bande.
\sfield{Coûts} Loss-versus-rebalancing (26\,000~\$ par semaine pour 20~millions de dollars avec une volatilité quotidienne de 4~\%), couverture, gas.
\sfield{Comment elle s'éteint} Les pics de volatilité, qui font croître la perte avec le carré de la volatilité, alors que les frais suivent le volume.
\sfield{Horizon, capacité, infrastructure} De quelques jours à quelques mois.
\sfield{Backtest honnête} Volume de bruit qui réagit au palier de frais.
\sfield{Sources} Milionis, Moallemi, Roughgarden et Zhang.
\end{strategyfile}

\section{Tutoriel~: céder la valeur au constructeur par enchère}\label{tut:hf:on-chain-trading:tut}

\begin{tutorial}
\textbf{Objectif.} Valoriser l'arbitrage, l'\omterm{def:hf:on-chain-trading:bid}{enchère de bundle} et la semaine du fournisseur de liquidité couvert sur une chaîne simulée. \textbf{État final~:} les
figures et les nombres du texte.
\begin{tutsteps}
\item \textbf{L'arbitrage optimal} en forme fermée, vérifié contre le swap entier de \texttt{firm.amm}.
\omcode{code/firm/searcher/firm_searcher.py}{38}{51}{Vendre dans un pool cher ou acheter dans un pool bon marché, jusqu'à ce que l'unité marginale rapporte le prix centralisé.}\label{lst:hf:on-chain-trading:arb}
\item \textbf{L'enchère} entre $n$ searchers.
\omcode{code/firm/searcher/firm_searcher.py}{66}{75}{Valeurs privées uniformes et enchère d'équilibre au premier prix~; la part du constructeur.}\label{lst:hf:on-chain-trading:auction}
\item \textbf{La semaine du pool}~: swaps de bruit, arbitrage après chaque bloc, le compte du fournisseur couvert.
\omcode{code/firm/searcher/firm_searcher.py}{103}{131}{Swap de bruit et arbitrage de chaque bloc, valorisés au prix centralisé.}\label{lst:hf:on-chain-trading:pool}
\end{tutsteps}
\textbf{Pour aller plus loin.} Laisser le volume de bruit baisser avec les frais~; ajouter un second pool de liquidité et chercher les cycles atomiques~; intégrer dans leurs coûts la latence vers la
plateforme centralisée (CEX).
\end{tutorial}

\section{Construction~: le searcher}\label{bld:hf:on-chain-trading:searcher}

\begin{build}
\bfield{Rôle} Valoriser l'arbitrage on-chain et les liquidations, l'\omterm{def:hf:on-chain-trading:bid}{enchère de bundle} qui les attribue et la perte qu'y subit un fournisseur couvert.
\bfield{Interface} \texttt{optimal\_arb(x, y, p, fee)}, \texttt{arb\_exact}, \texttt{auction(value, n, spread, seed, trials)},
\texttt{Chain(blocks, seed, sigma\_day, p0, noise\_usd)}, \texttt{run\_pool(chain, value0, fee, gas\_usd)}, \texttt{lvr\_rate}, \texttt{liquidation\_value},
\texttt{gas\_cost}. Repose sur \texttt{firm.amm} et \texttt{firm.gasfee} du livre~3~; l'arithmétique de l'attaque sandwich relève de \texttt{firm.sandwich}, également du livre~3.
\bfield{Règles} Dollars~; frais exprimés en fractions~; un arbitrage à chaque bloc, après le swap de bruit.
\bfield{Tests d'acceptation} \texttt{code/firm/searcher/tests/}~: la condition du premier ordre de l'arbitrage et son optimalité face aux tailles voisines~; aucun
arbitrage dans la bande de frais~; accord avec le swap entier~; la répartition de l'enchère à la main~; le loss-versus-rebalancing à 40~\% près de la formule sur
trois jours~; le compte couvert est cohérent.
\bfield{Pour aller plus loin} Cycles atomiques~; enchères de flux d'ordres~; coûts de searcher dépendant de la latence.
\end{build}

\begin{omsources}
\item P.~Daian et al., Flash Boys 2.0: frontrunning, transaction reordering, and consensus instability in decentralized exchanges,
arXiv:1904.05234, 2019 (IEEE Symposium on Security and Privacy, 2020).
\item J.~Milionis, C.~C.~Moallemi, T.~Roughgarden, A.~L.~Zhang, Automated market making and loss-versus-rebalancing, arXiv:2208.06046.
\item The Block, compte rendu du procès nul dans \emph{United States v.\ Peraire-Bueno}, 8 novembre 2025.
\end{omsources}

\section{Exercices}

\begin{exercise}[$\star$]\label{exo:hf:on-chain-trading:1}
Un pool détient 1\,000 jetons et 3\,100\,000~\$~; le prix centralisé est de 3\,000~\$ et les frais de 0.3~\%. Combien de jetons l'arbitragiste doit-il vendre dans
le pool, et pour quel profit~?
\end{exercise}

\begin{exercise}[$\star$]\label{exo:hf:on-chain-trading:2}
Quel est le taux de loss-versus-rebalancing d'un pool à produit constant de 20~millions de dollars avec une volatilité quotidienne de 4~\%, par jour et par semaine~?
\end{exercise}

\begin{exercise}[$\star$]\label{exo:hf:on-chain-trading:3}
Liquider 1~million de dollars de dette avec un bonus de 5~\% coûte 20~\$ de gas. Quelle est sa valeur~?
\end{exercise}

\begin{exercise}[$\star\star$]\label{exo:hf:on-chain-trading:4}
Expliquer pourquoi la part du constructeur tend vers 100~\% à mesure qu'on ajoute des searchers.
\end{exercise}

\begin{exercise}[$\star\star$]\label{exo:hf:on-chain-trading:5}
Expliquer pourquoi le loss-versus-rebalancing dépend à peine des frais du pool dans le modèle.
\end{exercise}

\begin{exercise}[$\star\star$]\label{exo:hf:on-chain-trading:6}
Expliquer pourquoi un arbitrage CEX--DEX n'est pas atomique et quel risque il laisse subsister.
\end{exercise}

\begin{exercise}[$\star\star\star$]\label{exo:hf:on-chain-trading:7}
\emph{Programmation.} Doubler la volatilité du jeton dans \texttt{firm.searcher.Chain}, puis expliquer l'effet sur le loss-versus-rebalancing et sur le résultat net du fournisseur couvert.
\end{exercise}

\begin{exercise}[$\star\star\star$]\label{exo:hf:on-chain-trading:8}
\emph{Trouver la faille.} «~Les opportunités d'arbitrage sur nos pools valaient 11\,400~\$ la semaine dernière~; notre searcher en fera autant.~»
\end{exercise}

\section{Problème~: céder la valeur au constructeur}

\begin{problem}[{Problème du week-end --- céder la valeur au constructeur}]\label{pb:hf:on-chain-trading:1}
Une entreprise fait tourner des searchers et une position d'\omterm{def:hf:on-chain-trading:hedged}{apport de liquidité couvert} sur une plateforme décentralisée (DEX).

\medskip\noindent\textbf{Partie I --- Arbitrage.}
\begin{enumerate}
\item Dériver la taille optimale de transaction CEX--DEX.
\item Définir l'\omterm{def:hf:on-chain-trading:atomic}{arbitrage atomique} et le comparer à l'arbitrage CEX--DEX.
\item Qu'ont documenté Daian et ses coauteurs~?
\item Décrire le chemin d'un bundle, du searcher jusqu'au bloc.
\end{enumerate}

\medskip\noindent\textbf{Partie II --- Enchères.}
\begin{enumerate}[resume]
\item Définir une \omterm{def:hf:on-chain-trading:bid}{enchère de bundle} (bundle bid).
\item Donner l'enchère d'équilibre avec des valeurs privées uniformes, et la part du constructeur pour 2, 5 et 20 searchers.
\item Que garde le gagnant, et pourquoi~?
\item Comment une liquidation est-elle mise aux enchères~?
\end{enumerate}

\medskip\noindent\textbf{Partie III --- Liquidité.}
\begin{enumerate}[resume]
\item Définir l'\omterm{def:hf:on-chain-trading:hedged}{apport de liquidité couvert} (hedged liquidity provision) et le taux de loss-versus-rebalancing.
\item Comparer la perte simulée à la formule.
\item Que prennent encore les arbitragistes au pool, et qui en a payé le prix~?
\item Pourquoi la comparaison des frais du modèle est-elle flatteuse pour les frais élevés~?
\end{enumerate}

\medskip\noindent\textbf{Partie IV --- Le verdict.}
\begin{enumerate}[resume]
\item Énoncer le \emph{résultat nommé}~: la part du profit d'arbitrage versée aux constructeurs à mesure que le nombre de searchers concurrents augmente, et le
rendement net du fournisseur de liquidité couvert après le loss-versus-rebalancing.
\item Que faut-il à un searcher pour garder plus de 1~\%~?
\item Que montre l'exemple de l'attaque sandwich sur qui paie qui~?
\item Résumer l'encadré daté.
\item Quelle fiche de stratégie est la plus exposée au flux d'ordres privé~?
\item Comment choisir le palier de frais d'un pool~?
\item Que changerait une enchère de flux d'ordres pour le swapper~?
\item En une phrase~: qui touche les gains d'un arbitrage on-chain~?
\end{enumerate}
\end{problem}

\section{Questions d'entretien}

\begin{interviewq}[$\star$ \iqroles{trader}]\label{iq:hf:on-chain-trading:1}
L'ether bouge de 1~\% sur les plateformes centralisées. Que se passe-t-il dans un pool à produit constant d'ether au bloc suivant, et qui en profite~?
\end{interviewq}

\begin{interviewq}[$\star\star$ \iqroles{chercheur}]\label{iq:hf:on-chain-trading:2}
Dérivez le taux de loss-versus-rebalancing d'un pool à produit constant.
\end{interviewq}

\begin{interviewq}[$\star\star$ \iqroles{développeur}]\label{iq:hf:on-chain-trading:3}
Concevez un searcher qui repère, simule et enchérit sur des bundles d'arbitrage dans le temps d'un bloc.
\end{interviewq}

\begin{interviewq}[$\star\star$ \iqroles{risque}]\label{iq:hf:on-chain-trading:4}
La jambe on-chain de votre arbitrage CEX--DEX n'a pas été incluse, mais la couverture centralisée a été exécutée. Quelle est votre position, et que devrait faire le système~?
\end{interviewq}

\begin{interviewq}[$\star\star$ \iqroles{trader}]\label{iq:hf:on-chain-trading:5}
Comment enchéririez-vous face à un concurrent dont les coûts sont inférieurs aux vôtres~?
\end{interviewq}

\begin{interviewq}[$\star\star\star$ \iqroles{chercheur}]\label{iq:hf:on-chain-trading:6}
Montrez que, avec $n$ enchérisseurs dont les valeurs sont uniformes sur $[a,b]$, enchérir $a+\frac{n-1}{n}(v-a)$ est un équilibre symétrique de l'enchère au
premier prix.
\end{interviewq}
